\section{A proved endoscopic injection and the origin of the small eigenvalue}
\label{sec:kt-endoscopic}

This section uses the same standard compact unitary group and full
lattice stabilizer as KT-U3-DIM. Write $\epsilon=\chi_{-3}$,
$E=\mathbf Q(\omega)$, and let $e\geq2$ be even.
For a normalized primitive form
$f\in S_{2e+1}(\Gamma_0(3),\epsilon)$, write $\pi_f$ for its
unitary automorphic representation of $\mathrm{GL}_2(\mathbf A_{\mathbf Q})$.
Thus its central character is $\epsilon$, and at a good prime the
sum of its unitary Satake roots is $a_p(f)/p^e$.

\begin{theorem}[KT-U3-ENDO-INJECT]
Primitive automorphic twisting by $\epsilon$ is an involution on
the normalized newforms of $S_{2e+1}(\Gamma_0(3),\epsilon)$.
Its fixed points are exactly the CM forms. For each two-element
orbit $\{f,f\otimes\epsilon\}$ there is a Hecke eigenline in
$M_{e-1}$ with generic endoscopic parameter
\[
             \psi_f=1\boxplus\operatorname{BC}_{E/\mathbf Q}(\pi_f).
\]
Different two-element orbits give independent eigenlines.
At every split prime $p\ne3$ their scaled eigenvalue is
\[
                  \alpha_p=p^e+a_p(f).                 \tag{KT-E.1}
\]
Consequently the dimension of the endoscopic part of $M_{e-1}$
is at least
\[
 \frac{r_e-c_e}{2}=\left\lfloor\frac{e-1}{3}\right\rfloor,
 \qquad
 r_e=\left\lfloor\frac{2e-2}{3}\right\rfloor,
 \quad c_e=\mathbf1_{3\mid e}.                         \tag{KT-E.2}
\]
For $e=4$ and $e=6$, this is the entire one-dimensional compact
space. In particular its stable part is zero in those two degrees.
\end{theorem}

\begin{proof}
\paragraph{Twists, base change, and the local level.}
The exact cyclic base-change inputs are Langlands,
\emph{Base change for $\mathrm{GL}(2)$}, author retypeset version,
Section~2, properties (A)--(F), p.~13, especially (C), and
Lemma~11.3(b), p.~133, Proposition~11.4, p.~136, and
Lemma~11.7, p.~142. They give existence with local compatibility,
cuspidality outside the dihedral exception, the two-element
quadratic-twist fibres, and the characterization of a nontrivial
quadratic self-twist by automorphic induction from that quadratic
field. These statements were read directly; the entire trace-formula
proof of base change was not independently reconstructed here.

The level and primitive nebentypus force the local conductor of
$\pi_f$ at three to be one and its determinant on inertia to be
the ramified quadratic character. Its Weil--Deligne representation
there has zero monodromy and inertial characters $1,\epsilon$.
Here is the elementary conductor check: the nontrivial determinant
forces the Artin conductor of the Weil representation to be at
least one. The Weil--Deligne conductor is this Artin conductor
plus the nonnegative monodromy term, so conductor one forces
Artin conductor one, one inertia invariant, and no Swan
contribution. Semisimple tame inertia is therefore exactly
$1\oplus\epsilon$. Monodromy commutes with inertia; on these
two distinct one-dimensional inertia lines a nilpotent commuting
operator is zero. Frobenius preserves the two distinct inertia lines,
so the local parameter is a principal-series pair
$\mu\oplus\nu\epsilon$ with $\mu,\nu$ unramified.
The conductor and local-parameter identification use the same
classical newform/local-compatibility input as K17-GL2-CONDUCTOR.

Twisting this pair by $\epsilon$ again gives conductor one.
At other finite primes the twist is unramified, and its weight
and nebentypus are unchanged. Hence the associated primitive
normalized twist has level three. This is an assertion about
the primitive automorphic twist; the coefficientwise Dirichlet
twist at multiples of three need not itself be that primitive
newform.
Langlands's self-twist criterion and KT-GL2-CM now show that
the fixed points are precisely the CM forms, all of which here
are induced from $E$.
For a nonfixed orbit the base change $\Pi=\operatorname{BC}(\pi_f)$
is cuspidal. Its finite local parameters are unramified:
the only ramified inertial character just exhibited is killed
on restriction to $W_{E_3}$, and there is no monodromy.
Its central character is $\epsilon\circ N_{E/\mathbf Q}=1$.
Also $\Pi^c\simeq\Pi$ by base change and
$\Pi^\vee\simeq\Pi$ by the rank-two trivial determinant,
so $\Pi$ is conjugate self-dual.
Its infinity characters are $(z/\bar z)^e$ and $(z/\bar z)^{-e}$,
distinct positive-parity characters. Here the global sign deduction
uses both KMSW seed theorems: Proposition~1.3.1, p.~67,
attaches $\Pi$ to the unique embedding $\eta_{\chi_\kappa}$,
and Proposition~1.3.3, p.~71, localizes this same embedding
at infinity. The local image has parity $(-1)^{2-1}\kappa$
by the bilinear-form definition in Section~1.2.3, p.~60,
and the construction in Section~1.2.4, Lemma~1.2.5, p.~61.
The two distinct one-dimensional positive-parity summands
cannot support a negative-parity form: such a form must be
diagonal on these distinct conjugate-self-dual characters, and
an alternating form on a one-dimensional space is zero.
Thus $-\kappa=1$. Proposition~1.3.1 identifies this as
global conjugate-orthogonal parity, so $\Pi$ is a permitted
positive-parity simple constituent of the rank-three parameter.
No even-rank application of Shin is needed.

The finite local parameters of $\Pi$ are tempered. At good
primes this follows from the classical Ramanujan bound.
At three the two principal-series characters of $\pi_f$ have
distinct restrictions to inertia. A unitary principal series
with such characters is tempered: the non-tempered unitary
principal series has inducing characters with the same unitary
part, and hence the same restriction to inertia. Base change
preserves this temperedness. Thus at every finite nonsplit
place the roots of $\Pi$ are $\{b,b^{-1}\}$, with $|b|=1$.

\paragraph{The compact packet member and its multiplicity.}
The global parameter $\psi=1\boxplus\Pi$ is discrete and generic.
Its reduced component group is of order two, generated by the
class of $a$, where $a$ acts as $-I_2$ on $\Pi$ and as $1$ on
the one-dimensional summand. At split finite places choose the
usual spherical representation. At inert places choose the
hyperspecial member of the unramified packet. At the ramified
place three, the full-stabilizer principal-series argument in
the coverage section supplies a $K_3$-fixed packet member for
the parameter $\{b,1,b^{-1}\}$. More explicitly, take the
unramified $E_3^\times$ character $\lambda_1$ with value $b$
and take $\lambda_2=1$ in Adler--Lansky's notation (4.4.1).
Section~4.4, pp.~14--15, excludes (4.4PS--1) by unitarity
and (4.4PS--2) because its restriction to $\mathbf Q_3^\times$
is ramified. The only reducible possibility is (4.4PS--3),
the nontrivial unramified quadratic character $b=-1$.
For $b=1$ the principal series is irreducible.
Proposition~6.1(1),(4), p.~21, gives the one-member or
two-member packet, and Proposition~7.1(i),(iv), p.~28,
gives its base change. The induced $\mathrm{GL}_3$ representation
is irreducible even when roots repeat: a ratio of its unitary
characters cannot equal $|\cdot|^{\pm1}$.
Iwasawa decomposition for the full stabilizer and the triviality
of $\lambda$ on $B\cap K_3$ give one fixed vector in the
principal series; exactness puts it in a constituent of this
packet. These are local statements, without a global
cuspidality requirement on a rank-three base change.

The value of every finite local packet character on $a$ is one.
At a nonsplit place this can be checked directly in the
centralizer, without selecting a sign for the ramified
constituent. If $b\ne b^{-1}$, the pair $b,b^{-1}$ contributes
the connected torus $\{\operatorname{diag}(t,t^{-1})\}$,
which contains $-I_2$. If $b=b^{-1}$, the multiplicity-two
space has an orthogonal or symplectic form. The element
$-I_2$ lies respectively in $\mathrm{SO}_2(\mathbf C)$ or
$\mathrm{Sp}_2(\mathbf C)$, both connected. If it merges with
the trivial third character, $a$ has determinant one in
$O_3(\mathbf C)$ and lies in $\mathrm{SO}_3(\mathbf C)$.
All these connected subgroups lie in $\mathrm{SL}_3$.
The exact definition is KMSW Section~0.4.4, p.~27,
repeated in Section~1.6.1, p.~80:
\[
 S_\phi^{\mathrm{rad}}=
       \operatorname{Cent}(\phi,\widehat G_{\mathrm{der}})^\circ,
 \qquad S_\phi^{\natural}=S_\phi/S_\phi^{\mathrm{rad}}.
\]
Here $\widehat G_{\mathrm{der}}=\mathrm{SL}_3$.
Consequently $a$ belongs to $S_{\psi_v}^{\mathrm{rad}}$, so
it is trivial even in the group $S_{\psi_v}^{\natural}$ on
which Theorem~1.6.1(2), p.~80, places the packet characters.
At a split place the
spherical packet character has the same trivial value.

At infinity the ordered weights are $(e,0,-e)$ and the compact
member has coefficient highest weight $(e-1,0,1-e)$.
The real discrete-series packet rule says that for a compact
root between positions $i<j$, the packet character takes product
$(-1)^{i+j}$ on the two corresponding sign generators.
This is Gan--Gross--Prasad, \emph{Restrictions of representations
of classical groups: examples}, Ast\'erisque 346 (2012),
Section~2, p.~115, the compact-root formulas following
Corollary~2.3. In the compact group all roots are compact.
Only the known real local-packet description is used here,
not the restriction conjecture studied elsewhere in that paper.
The two characters from $\Pi$ occupy positions one and three,
so the value on $a$ is $(-1)^{1+3}=1$.
This product rule does not require choosing the individual signs
of the two generators.

The product of all local characters is therefore trivial on
the reduced global component group. KMSW Section~1.7,
p.~89, states that $\epsilon_\psi=1$ for a generic parameter.
The proved generic/pure-inner case of Theorem~5.0.5, p.~206,
therefore places the chosen tensor product in the discrete
automorphic spectrum with multiplicity one. Its finite
$K$-fixed vector and its compact coefficient type give a
Hecke eigenline in $M_{e-1}$.

\paragraph{Independence, counting, and the eigenvalues.}
Langlands's fibre assertion (C) identifies two cuspidal base
changes exactly when their original forms differ by the
quadratic twist. Distinct two-element orbits therefore give
distinct global parameters. KMSW puts these in distinct
direct summands, so the selected eigenlines are independent.
The count of orbits is $(r_e-c_e)/2$ by KT-GL2-CM;
substitution in the three even residue classes modulo six
gives the floor in (KT-E.2).

At a split prime the three unitary roots consist of $1$ and
the two unitary roots of $\pi_f$, so their sum is
$1+a_p(f)/p^e$. Multiplication by $p^e$ in the already checked
compact Hecke normalization gives (KT-E.1).
Finally KT-U3-DIM gives $D(3)=D(5)=1$, while (KT-E.2)
gives one line at $e=4,6$. Hence these full spaces are
endoscopic, and neither contains a stable summand.
\end{proof}

\begin{corollary}[KT-U3-ENDO-SIX]
The unique compact form with Harish--Chandra parameter
$(6,0,-6)$, equivalently coefficient highest weight
$(5,0,-5)$, is the endoscopic lift of the non-CM quadratic-twist
pair in $S_{13}(\Gamma_0(3),\chi_{-3})$.
It has
\[
  7^5T_7=7^6+40250=157899,
  \qquad
  13^5T_{13}=13^6+1284050=6110859.
\]
\end{corollary}
\begin{proof}
The classical space has dimension three and CM dimension one.
The exact classical calculation K17-GL2-MODEL, independently
recomputed in the report's classical section, gives $a_7=40250$
on its two non-CM eigenforms; the CM eigenvalue is $-153502$.
The corresponding non-CM value at thirteen is $1284050$.
KT-U3-ENDO-INJECT and the one-dimensional compact dimension
identify the complete global parameter and give both displayed
values. This argument does not infer endoscopy from agreement
at finitely many primes; it first constructs the global
endoscopic eigenline and then uses the exact dimension.
\end{proof}

\paragraph{The all-weight converse.}
The injection proves the lower bound (KT-E.2), and already
proves the entire stable/endoscopic split at $e=4,6$.
The following section supplies the converse classification,
including the outer-character and other parameter cases.
KT-U3-ENDO-EXACT proves that this family exhausts the endoscopic
spectrum for every even $e\ge2$ and identifies
$D(e-1)-\lfloor(e-1)/3\rfloor$ as the stable dimension.
